\section{Polygons, chirotopes, chambers, and walls}\label{sec:polygons}

\subsection{Conventions}\label{sec:conventions}

For $u=(u_1,u_2)$, $v=(v_1,v_2)\in\RR^2$ put $\det(u,v)=u_1v_2-u_2v_1$.
The sign function is $\sgn\colon\RR\to\{-1,0,1\}$. The step function is
\[
\Theta\colon\RR\setminus\{0\}\to\{0,1\},\qquad
\Theta(x)=1\ (x>0),\qquad \Theta(x)=0\ (x<0);
\]
it is never evaluated at $0$ in this document, and every statement that
consults it is made on a domain where its argument is nonzero
(Lemma~\ref{lem:gates-nonzero}). For $n\geq3$ indices of vertices and edges
are elements of $\ZZ/n$; $i+1$, $i-1$ are read modulo $n$. The main text
writes $[u,v]=u_2v_1-u_1v_2=-\det(u,v)$; this document uses $\det$ only.
The frozen source CV uses one-based tail labels: its edge $e_i$ runs
from $p_i$ to $p_{i+1}$, so $p_i=\mu_i$ requires no index shift.
RC instead uses zero-based head labels: $e_r$ runs from $p_{r-1}$
to $p_r$, and its starred bracket is $[u,v]_*=-\det(u,v)$.
Under $p_r=\mu_{r+1}$ an RC vertex $r$ becomes vertex $r+1$;
its incoming edge $e_r$ is the SM tail edge $E_r$, with residue zero
represented by $n$. Here $E_i=[\mu_i,\mu_{i+1}]$ and
$\lt_i=\mu_{i+1}-\mu_i$, as in the main text.

\subsection{Polygons and chirotopes}

\begin{definition}[labelled tuples and polygons]\label{def:polygon}
Let $n\geq3$. A \emph{labelled tuple} is an element
$P=(\mu_1,\ldots,\mu_n)\in(\RR^2)^n$, indices in $\ZZ/n$ with $\mu_{n+1}=\mu_1$. The
\emph{cyclic shift} is $(\sigma P)_i=\mu_{i+1}$. A \emph{polygon} on $n$
vertices is an orbit $[P]=\{\sigma^kP:k\in\ZZ/n\}$ of labelled tuples; a
labelled tuple in the orbit is a \emph{representative}, and the space of
polygons is the quotient $(\RR^2)^n/(\ZZ/n)$. Every notion below is
introduced on a representative and is invariant under $\sigma$ (up to the
induced relabelling of indices, which is recorded where it matters); it
therefore defines a notion on polygons. A polygon has no distinguished vertex
and no distinguished edge; when a construction needs one (a root, a base
vertex), that choice is part of the construction's data.

For a representative, the \emph{edge vectors} are
\[
\lt_i=\mu_{i+1}-\mu_i\qquad(i\in\ZZ/n),
\]
so that $\sum_i\lt_i=0$, and the \emph{edge segments} are
$E_i=[\mu_i,\mu_{i+1}]=\{\mu_i+t\lt_i:0\leq t\leq1\}$. Vertex $i$ is
\emph{incident} to $E_{i-1}$ and $E_i$. Two edges $E_i,E_j$ are
\emph{adjacent} if $j-i\in\{-1,0,1\}$ and \emph{remote} otherwise; an edge
$E_j$ is \emph{remote to the vertex} $i$ if $j\notin\{i-1,i\}$, and
\emph{remote to the vertex and its two edges} if $j\notin\{i-2,i-1,i,i+1\}$.
No condition is imposed on the tuple by this definition; in particular
coincident vertices and zero edges are allowed at this level.
\end{definition}

\begin{remark}[formalization]
The quotient by the finite group $\ZZ/n$ is the only quotient used. On the
generic locus the action is free (distinct vertices), so the quotient is a
$2n$-manifold and its connected components are the orbits of the components
of the labelled generic locus. A formalization may work with labelled tuples
throughout and record each construction's shift-equivariance as a lemma; the
statements below are written so that both readings agree.
\end{remark}

\begin{definition}[chirotope and turns]\label{def:chirotope}
For $i,j,k\in\ZZ/n$ put
\[
\chi_{ijk}(P)=\sgn\det(\mu_j-\mu_i,\ \mu_k-\mu_i)\in\{-1,0,1\}.
\]
The \emph{turn} at vertex $i$ is $\tau_i(P)=\chi_{i-1,i,i+1}(P)$. The number
of \emph{left turns} is $\ell(P)=\#\{i:\tau_i(P)=1\}$.
\end{definition}

\begin{lemma}[elementary properties of $\chi$]\label{lem:chi-basic}
For all $i,j,k$:
\begin{enumerate}
\item[(i)] $\chi_{ijk}$ is alternating: it changes sign under a transposition
of two of its indices and is unchanged under a cyclic permutation; it is $0$
when two indices coincide.
\item[(ii)] $\chi_{ijk}=\sgn\bigl([\mu_j-\mu_i,\ \mu_j-\mu_k]\bigr)$ with the
main text's bracket $[u,v]=u_2v_1-u_1v_2$.
\item[(iii)] $\chi_{i,i+1,k}=\sgn\det(\lt_i,\ \mu_k-\mu_i)$; when $\lt_i\neq0$
we say (a naming convention) that $\mu_k$ lies \emph{strictly to the left} of
the directed line through $E_i$ if this sign is $+1$ and \emph{strictly to the
right} if it is $-1$.
\item[(iv)] $\tau_i=\sgn\det(\lt_{i-1},\lt_i)$; $\tau_i=+1$ is a left turn.
\end{enumerate}
\status{proved (refereed: bench A 2026-09-12T06:44:52Z, bench B 2026-09-12T06:57:24Z; SM1-carried, read on SM12, repaired in round 12 (F-25-194), re-read on SM13)}
\end{lemma}
\begin{proof}
(i) $\det$ is alternating and $\det(\mu_j-\mu_i,\mu_k-\mu_i)$ is the standard
area form of the ordered triple; a cyclic permutation of the triple is even.
(ii) $[\mu_j-\mu_i,\mu_j-\mu_k]=-\det(\mu_j-\mu_i,\mu_j-\mu_k)
=\det(\mu_j-\mu_i,\mu_k-\mu_j)=\det(\mu_j-\mu_i,\mu_k-\mu_i)$, the last step
by $\det(u,v-u)=\det(u,v)$. (iii) is the case $j=i+1$ of the definition.
(iv) $\det(\mu_i-\mu_{i-1},\mu_{i+1}-\mu_{i-1})=\det(\lt_{i-1},\lt_{i-1}+\lt_i)
=\det(\lt_{i-1},\lt_i)$.
\end{proof}

\subsection{Generic polygons and chambers}

\begin{definition}[generic polygon]\label{def:generic}
A polygon $P$ is \emph{generic} if
\begin{enumerate}
\item[(G1)] $\chi_{ijk}(P)\neq0$ for all pairwise distinct $i,j,k\in\ZZ/n$;
\item[(G2)] there is no point of $\RR^2$ lying in the relative interiors of
three distinct edge segments.
\end{enumerate}
The set of generic polygons on $n$ vertices is $\mathcal U_n\subseteq(\RR^2)^n$.
\end{definition}

\begin{lemma}[what (G1) gives]\label{lem:g1}
Let $P$ satisfy (G1). Then:
\begin{enumerate}
\item[(i)] all vertices are distinct and every edge vector is nonzero;
\item[(ii)] every turn $\tau_i$ is $\pm1$; consecutive edge vectors are
linearly independent, so no edge vector is a positive or negative multiple
(``flat'' or ``kink'') of its predecessor;
\item[(iii)] no vertex lies on the line spanned by an edge not incident to it;
in particular no vertex lies on a non-incident closed edge segment;
\item[(iv)] two distinct adjacent edges meet exactly in their common vertex;
\item[(v)] two remote edges are either disjoint or meet in exactly one point,
which lies in the relative interior of both and at which the two edge
directions are linearly independent.
\end{enumerate}
\status{proved (refereed: bench A 2026-09-12T06:44:52Z, bench B 2026-09-12T06:57:24Z; SM1-carried, read on SM12, repaired in round 12 (F-25-190, F-25-195), re-read on SM13)}
\end{lemma}
\begin{proof}
(i) If $\mu_i=\mu_j$ with $i\neq j$, then $\chi_{ijk}=0$ for any third index
$k$ (there is one since $n\geq3$). (ii) $\tau_i=\sgn\det(\lt_{i-1},\lt_i)\neq0$.
(iii) If $\mu_k$ lies on the line through $\mu_i,\mu_{i+1}$ with
$k\notin\{i,i+1\}$ then $\chi_{i,i+1,k}=0$. (iv) By (ii) the two edges at a
vertex span the plane, so their lines meet only at the vertex. (v) Let
$E_i,E_j$ be remote and $x\in E_i\cap E_j$. By (iii) $x$ is not an endpoint
of either. If $\lt_i,\lt_j$ were parallel, the two segments would lie on one
line, and then an endpoint of one would lie on the line of the other, contrary
to (iii). Two segments on non-parallel lines meet in at most one point.
\end{proof}

\begin{definition}[crossings]\label{def:crossings}
Let $P$ satisfy (G1). The \emph{crossing set} $X(P)$ is the set of unordered
pairs $\{i,j\}$ of remote edge indices with $E_i\cap E_j\neq\varnothing$. For
$\{i,j\}\in X(P)$ the \emph{crossing point} $x_{ij}$ is the unique point of
$E_i\cap E_j$, and its \emph{parameters} $t_{ij}\in(0,1)$ on $E_i$ and
$t_{ji}\in(0,1)$ on $E_j$ are defined by $x_{ij}=\mu_i+t_{ij}\lt_i
=\mu_j+t_{ji}\lt_j$. The \emph{crossing sign} of $\{i,j\}$ read from $E_i$ is
$\sgn\det(\lt_i,\lt_j)$.
\end{definition}

\begin{lemma}[crossings through the chirotope]\label{lem:crossing-test}
Let $P$ satisfy (G1) and let $E_i,E_j$ be remote. Then $\{i,j\}\in X(P)$ if
and only if
\[
\chi_{i,i+1,j}\,\chi_{i,i+1,j+1}=-1\quad\text{and}\quad
\chi_{j,j+1,i}\,\chi_{j,j+1,i+1}=-1.
\]
If $P$ is generic, distinct crossings have distinct crossing points, and
$X(P)$ is finite. \status{proved (refereed: bench A 2026-09-12T06:44:52Z, bench B 2026-09-12T06:57:24Z; SM1-carried, read on SM12, repaired in round 12 (F-25-196), re-read on SM13)}
\end{lemma}
\begin{proof}
By Lemma~\ref{lem:g1}(i),(iii) the edge vectors $\lt_i,\lt_j$ are nonzero and
none of the four endpoints of $E_i,E_j$ lies on the line of the other segment,
so the four chirotope entries in the display are $\pm1$. Separation
criterion: along $E_j$ the signed distance from the line of $E_i$,
$f(t)=\det(\lt_i,\ \mu_j+t\lt_j-\mu_i)$, is affine in $t$ with
$f(0),f(1)\neq0$, so $E_j$ meets the line of $E_i$ exactly when
$f(0)f(1)<0$, that is, when $\chi_{i,i+1,j}\,\chi_{i,i+1,j+1}=-1$
(Lemma~\ref{lem:chi-basic}(iii)); then $f$ has nonzero slope
$\det(\lt_i,\lt_j)$, the two lines meet in one point, and that point lies in
the relative interior of $E_j$. Symmetrically, $E_i$ meets the line of $E_j$
exactly when the second condition holds. If both conditions hold, the common
point of the two lines lies in the relative interior of both segments, so
$E_i\cap E_j\neq\varnothing$; conversely, a point of $E_i\cap E_j$ is a point
of $E_j$ on the line of $E_i$ and a point of $E_i$ on the line of $E_j$, so
both conditions hold. If two distinct crossings had the same point, that
point would lie in the relative interior of three distinct edges (the two
crossings share at most one edge), contrary to (G2). Finiteness is clear.
\end{proof}

\begin{lemma}[stability of separated segment intersections]
\label{lem:wall-segment-stability}
Consider finitely many continuously varying nonzero oriented segments.
For a specified pair, suppose that at the centre the segments are either
disjoint or meet transversely in both relative interiors. In the first case
they stay disjoint nearby. In the second case their unique intersection
persists, remains transverse and interior, and its two parameters vary
continuously. For finitely many such persistent crossings on a segment,
every pair with distinct central parameters keeps its parameter order nearby.
\status{new (round~4: no frozen-source statement carries this lemma --- the former locator RC tr:five-events is the five-event classification and does not state it, and the persistence argument appears only inside proofs (CV lem:dictionary, clause~(b); RC guard:canonical); the round-1 proof printed here is self-contained; F-25-139; cleared by bench~A on SM4 (2026-09-05T23:49:15Z))}
\end{lemma}
\begin{proof}
Disjoint compact segments have positive distance. Moving their endpoints by
at most $\eta$ moves every convex combination of the endpoints by at most
$\eta$, so a sufficiently small perturbation preserves their disjointness.
For the second case write the segments as $A+uD$ and $B+vH$, where
$0\leq u,v\leq1$. Solving $A+uD=B+vH$ by taking determinants gives
\begin{equation}\label{wallgeo:parameters}
 u=\frac{\det(B-A,H)}{\det(D,H)},\qquad
 v=\frac{\det(B-A,D)}{\det(D,H)}.
\end{equation}
The central denominator is nonzero. These quotients are therefore continuous
near the centre, and their central values in $(0,1)$ remain in that interval.
A nonzero continuous difference of two central parameters has constant sign
after shrinking. There are only finitely many requirements, so one common
interval satisfies all of them.
\end{proof}

\begin{definition}[chambers]\label{def:chamber}
A \emph{chamber} is a connected component of the space of generic polygons
$\mathcal U_n/(\ZZ/n)$; its preimage in $\mathcal U_n$ is a union of at most
$n$ components of $\mathcal U_n$ permuted by $\sigma$, the \emph{labelled
chambers}.
\end{definition}

\begin{proposition}[chambers]\label{prop:chambers}
$\mathcal U_n$ is open in $(\RR^2)^n$, and every chamber is open and path
connected. Along any path in $\mathcal U_n$ the following data are constant:
the chirotope $(\chi_{ijk})$; the crossing set $X(P)$; and, for every edge
$E_i$, the order of the parameters $t_{ij}$, $\{i,j\}\in X(P)$, along $E_i$.
\status{new (round~5: the two persistence clauses of the proof cited to Lemma~\ref{lem:wall-segment-stability}, F-25-143; the remainder is the SM1 proof, audited there)}
\end{proposition}
\begin{proof}
The finitely many determinants $\det(\mu_j-\mu_i,\mu_k-\mu_i)$ are continuous
and nonzero on $\mathcal U_n$, so (G1) is open. Condition (G2) is open as well:
if three edge interiors had a common point at a limit of polygons, the three
segments would meet in the limit; we argue directly. Fix $P\in\mathcal U_n$.
By Lemma~\ref{lem:g1}(v) each pair of remote edges is either disjoint or
meets transversally in one interior point; the first condition persists under
small perturbation, and in the second the crossing persists and varies
continuously (Lemma~\ref{lem:wall-segment-stability}).
The finitely many crossing points of $P$ are pairwise distinct
(Lemma~\ref{lem:crossing-test}), hence remain distinct nearby; a common point
of three edge interiors nearby would be a coincidence of two crossing points
(any two of the three edges are remote, since an adjacent pair meets only at
its vertex by Lemma~\ref{lem:g1}(iv), and a vertex on a remote edge is excluded
by (G1) nearby). So $\mathcal U_n$ is open. An open subset of $\RR^{2n}$ is
locally path connected, so its components are open and path connected.

Along a path in $\mathcal U_n$ each $\chi_{ijk}$ is a continuous nonzero
integer-valued function, hence constant. The crossing set is a function of the
chirotope (Lemma~\ref{lem:crossing-test}). Two crossings on a common edge $E_i$
change their order only if their parameters coincide at some time; at that
time the two crossing points coincide, which puts a common point in three edge
interiors, contrary to (G2) --- the two other edges are remote to $E_i$ and,
if they were adjacent to each other, their only common point is their shared
vertex, which cannot lie on $E_i$ by (G1).
\end{proof}

\subsection{The Gauss word, interlacement, and the visible signature}

\begin{definition}[traversal circle and Gauss word]\label{def:gauss}
Let $P$ be generic. The \emph{traversal circle} $\Gamma(P)$ is the set of
pairs $(i,t)$, $i\in\ZZ/n$, $t\in[0,1)$, with the cyclic order in which
$(i,t)<(i,t')$ for $t<t'$ and $(i,t)<(i+1,t')$ for all $t,t'$; the point of
the plane at $(i,t)$ is $\mu_i+t\lt_i$. Each crossing $\{i,j\}\in X(P)$ has
two \emph{visits} $(i,t_{ij})$ and $(j,t_{ji})$ on $\Gamma(P)$. The
\emph{Gauss word} of $P$ is the cyclic sequence of the $2|X(P)|$ visits in
the order of $\Gamma(P)$, each visit labelled by its crossing.
\end{definition}

\begin{definition}[interlacement]\label{def:interlace}
Two distinct crossings $x,y\in X(P)$ \emph{interlace} if the four visits
$x_0,y_0,x_1,y_1$ occur in the cyclic order $x_0<y_0<x_1<y_1$ (up to rotation),
that is, exactly one visit of $y$ lies between the two visits of $x$. The
\emph{interlacement graph} $G_P$ has vertex set $X(P)$ and an edge between
each interlacing pair. $\Ind(G_P)$ denotes the set of independent subsets of
$X(P)$, including $\varnothing$; for $S\subseteq X(P)$ let $N(S)$ be the set of
crossings interlacing some element of $S$, and let
$U(S)=X(P)\setminus(S\cup N(S))$.
\end{definition}

\begin{definition}[visible signature]\label{def:visible}
The \emph{visible signature} of a generic $P$ is the triple
\[
\bigl((\tau_i)_{i\in\ZZ/n},\ X(P),\ \text{the Gauss word of }P\bigr).
\]
By Proposition~\ref{prop:chambers} it is constant on chambers.
\end{definition}


\begin{definition}[weakly generic polygons and the silent locus]\label{def:weak}
A polygon $P$ is \emph{weakly generic} if every edge vector is nonzero, every
turn $\tau_i$ is nonzero, no vertex lies on a non-incident closed edge
segment, any two remote edges meet in at most one point and there
transversally, and (G2) holds. This is the main text's ``generic''. Every
generic polygon is weakly generic; the \emph{silent locus} is
$\mathcal S_n=\mathcal W_n\setminus\mathcal U_n$, where $\mathcal W_n$ is the
set of weakly generic polygons. A \emph{visible chamber} is a connected
component of $\mathcal W_n$.
\end{definition}

\begin{lemma}[openness of the weak locus]\label{lem:weak-open}
Read on labelled representatives, $\mathcal W_n$ is open in
$(\RR^2)^n$. Every connected component is open and polygonally
connected. In particular, every labelled weakly generic tuple has
an open ball contained in its visible chamber.
\status{new (round~4; adopted from the Codex proposal sm4\_round4\_core\_v1, P-25-23, F116-OPEN, re-read by the author; F-25-116)}
\end{lemma}
\begin{proof}
Fix a tuple $P$ satisfying Definition~\ref{def:weak}. Nonzero edge
vectors and nonzero turn determinants persist under small perturbations
by continuity. Thus distinct adjacent segments still meet only at their common
endpoint: their supporting lines are independent and have that point
in common.

Suppose every vertex moves by at most $\varepsilon$. A point of an
edge at affine parameter $t\in[0,1]$ moves by at most
$(1-t)\varepsilon+t\varepsilon=\varepsilon$, by the triangle
inequality. Consequently a vertex--segment distance, or a distance
between two segments, decreases by at most $2\varepsilon$.
Every vertex of $P$ has positive distance from every nonincident
closed segment, since the segment is compact and does not contain
the vertex. There are finitely many such pairs, so these separations
all persist nearby. In particular, remote segments cannot acquire
an endpoint contact.

Each disjoint pair of remote closed segments has positive compact
distance; preserve all these distances as well. This step allows
parallel or collinear disjoint segments. For an intersecting remote
pair, the endpoint exclusion puts the intersection in both relative
interiors, and weak genericity makes the two directions independent.
Write their initial points as $p,q$ and directions as $u,v$, and put
$a=q-p$ and $D=\det(u,v)\neq0$. The intersection equation is
$tu-sv=a$. Taking its determinant with $v$ gives
$tD=\det(a,v)$. Taking its determinant with $u$ gives
$sD=\det(a,u)$, since $-\det(v,u)=D$. Thus the two parameters are
\begin{equation}\label{eq:weak-open-parameters}
 t=\frac{\det(a,v)}D,\qquad s=\frac{\det(a,u)}D.
\end{equation}
The denominator stays nonzero nearby. Both parameters, and hence the
intersection point $p+tu$, vary continuously. Shrink the neighbourhood
so the parameters remain strictly between zero and one. The two
supporting lines then have exactly this one transverse intersection.

Intersection points belonging to distinct unordered pairs of remote edges at $P$ are
distinct: if two were equal, their union would contain at least three
distinct edges whose relative interiors contain that point, contrary
to (G2). Preserve the finitely many positive distances between these
points. Every nearby intersecting remote pair was already intersecting
at $P$, because the disjoint pairs remain disjoint. A nearby point in
three edge interiors would involve pairwise remote edges, since
adjacent interiors are disjoint. It would therefore identify the
points of two distinct preserved crossing pairs, a contradiction.
If there were initially fewer than two crossing pairs, the same
argument gives a contradiction before any separation needs to be
chosen. Thus (G2) persists.

Intersecting the finitely many neighbourhoods just constructed
preserves every condition of weak genericity. This proves openness.
A sufficiently small ball about any point is contained in the open
set and is convex; it therefore lies in the connected component of
its centre. Hence every component is open. Within one component,
fix a point and consider the points reachable from it by a finite
polygonal path. A small ball about a reachable point is reachable
by adding a final segment. A small ball about an unreachable point
cannot contain a reachable point, since a final segment would then
reach its centre. The reachable set and its complement are both
relatively open. Connectedness makes the reachable set the entire
component, proving polygonal connectedness.
\end{proof}

\begin{lemma}[the silent locus]\label{lem:silentlocus}
If $P$ is weakly generic and $\chi_{ijk}(P)=0$ for pairwise distinct
$i,j,k$, then either no two of $i,j,k$ are consecutive (a pure chord), or
exactly one pair is consecutive, say $\{i,i+1\}$, and $\mu_k$ lies on the
line of $E_i$ outside the closed segment $E_i$ (an exterior extension). The
visible signature (Definition~\ref{def:visible}) is defined and locally
constant on $\mathcal W_n$, hence constant on visible chambers, and
$\mathcal U_n$ is open and dense in $\mathcal W_n$. \status{new (round~5: the density clause cites Lemma~\ref{lem:weak-open} for the openness of $\mathcal W_n$ it relies on and names the polynomial whose zero set is avoided, F-25-147; the remainder is the SM1 proof, audited there)}
\end{lemma}
\begin{proof}
Two consecutive pairs in $\{i,j,k\}$ make it a turn triple, excluded. One
consecutive pair $\{i,i+1\}$ with $\mu_k$ on the line of $E_i$: $\mu_k$ is not
on the closed segment by hypothesis. Crossings are created or destroyed only
by a vertex passing through a remote closed segment, and their order along an
edge changes only through a coincidence of crossing points; both are excluded
on $\mathcal W_n$ exactly as in Proposition~\ref{prop:chambers}, so the
visible signature is locally constant. Density: $\mathcal W_n$ is open
(Lemma~\ref{lem:weak-open}) and $\mathcal U_n$ is its subset on which every
$\chi_{ijk}$ is nonzero; the zero set of the nonzero polynomial
$\prod\det(\mu_j-\mu_i,\mu_k-\mu_i)$ over the triples has empty interior in
$(\RR^2)^n$, so it contains no open subset of $\mathcal W_n$, and every point
of $\mathcal W_n$ is a limit of points of $\mathcal U_n$.
\end{proof}

\subsection{Rotation number}

\begin{definition}[regular locus and principal turns]\label{def:regular}
For a polygon $P$ and an index $i$ with $\lt_{i-1},\lt_i\neq0$ and $\lt_i$
not a negative real multiple of $\lt_{i-1}$, the \emph{principal turn}
$\vartheta_i(P)\in(-\pi,\pi)$ is the unique angle with
$\cos\vartheta_i=\langle\lt_{i-1},\lt_i\rangle/(|\lt_{i-1}||\lt_i|)$ and
$\sgn\vartheta_i=\sgn\det(\lt_{i-1},\lt_i)$. The \emph{regular locus}
$\mathcal R_n$ is the set of polygons with all $\lt_i\neq0$ and no $\lt_i$ a
negative multiple of $\lt_{i-1}$; equivalently, all principal turns exist.
Every generic polygon lies in $\mathcal R_n$ (Lemma~\ref{lem:g1}(ii)), and
$\vartheta_i(P)=0$ exactly when $\lt_i$ is a positive multiple of $\lt_{i-1}$.
\end{definition}

\begin{definition}[reversal]\label{def:shift}
For a representative $P=(\mu_1,\ldots,\mu_n)$ the \emph{reversal} is
$(\overline P)_i=\mu_{2-i}$ ($i\in\ZZ/n$), which fixes $\mu_1$ and reverses
the traversal. Reversal commutes with the shift up
to a shift, so it is defined on polygons; it is not a relabelling but a change
of orientation, and no law below is asserted to be invariant under it except
where stated. Write $\rho(P)=\overline P$ for this reversal map.
\end{definition}

\begin{lemma}[rotation number]\label{lem:rot}
For $P\in\mathcal R_n$ put $\rot(P)=\frac1{2\pi}\sum_{i\in\ZZ/n}\vartheta_i(P)$.
Then:
\begin{enumerate}
\item[(i)] $\rot(P)\in\ZZ$;
\item[(ii)] $\rot$ is constant along every continuous path in $\mathcal R_n$;
\item[(iii)] inserting a vertex in the relative interior of an edge does not
change $\rot$; reversing the traversal negates
it;
\item[(iv)] if $n=3$ then $\rot(P)=\tau_1(P)=\tau_2(P)=\tau_3(P)\in\{\pm1\}$
and in particular $\rot(P)\neq0$;
\item[(v)] $2|\rot(P)|<n$.
\end{enumerate}
\status{proved (refereed: bench A 2026-09-12T06:44:52Z, bench B 2026-09-12T06:57:24Z; SM1-carried, read on SM12, no text change, re-read on SM13)}
\end{lemma}
\begin{proof}
Identify $\RR^2$ with $\mathbb C$ and put $q_i=\lt_i/|\lt_i|$. The defining
property of the principal turn is $q_i=e^{\mathrm i\vartheta_i}q_{i-1}$.
Multiplying around the cycle gives $\exp(\mathrm i\sum_i\vartheta_i)=1$, so
the sum is in $2\pi\ZZ$: (i). The principal turn is a continuous function of
$(\lt_{i-1},\lt_i)$ on the open set where both are nonzero and not
antiparallel, so $\rot$ is a continuous integer-valued function on
$\mathcal R_n$: (ii). Inserting a vertex on an edge adds one zero turn and
changes no other turn; reversal reverses the cyclic order of the pairs and
replaces each $(\lt_{i-1},\lt_i)$ by $(-\lt_i,-\lt_{i-1})$, negating every
principal turn: (iii). For (iv), all three turn determinants of a triangle
equal $\det(\mu_2-\mu_1,\mu_3-\mu_1)$, which is nonzero for $P\in\mathcal R_3$
(three collinear points with nonzero edges have a doubled-back consecutive
pair); so the three turns have one sign, their sum lies strictly between $0$
and $3\pi$ in absolute value, and is a multiple of $2\pi$. For (v),
$2\pi|\rot(P)|\leq\sum_i|\vartheta_i|<n\pi$.
\end{proof}

\begin{lemma}[rotation of uniform and one-dissent polygons]\label{lem:uniformrot}
Let $P\in\mathcal R_n$ have all turns nonzero.
\begin{enumerate}
\item[(i)] If all turns are left, then $\rot(P)\geq1$; if all are right, then
$\rot(P)\leq-1$.
\item[(ii)] If exactly one turn is right and all others are left, then
$\rot(P)\geq1$; symmetrically with left and right exchanged, $\rot(P)\leq-1$.
\end{enumerate}
\status{proved (refereed: bench A, 2026-09-05T15:51:04Z; transcribed from CV lem:uniformrot (mirror clauses by reversal, C033))}
\end{lemma}
\begin{proof}
(i) All $\vartheta_i>0$ and there are at least three of them, so the sum is
positive and $\rot(P)$ is a positive integer. (ii) Let $\alpha\in(0,\pi)$ be
the magnitude of the right turn and $\Pi$ the sum of the left turns, so
$\Pi-\alpha=2\pi\rot(P)$. If $\rot(P)\leq0$ then $\Pi\leq\alpha<\pi$. Cut the
traversal after the right-turn vertex; the edge directions from there on are
obtained by successive positive turns summing to $\Pi<\pi$, so all edge
vectors lie in a closed angular sector of width $<\pi$, whose dual cone is
nonempty: there is $u$ with $\langle u,\lt_i\rangle>0$ for all $i$. Summing
contradicts $\sum_i\lt_i=0$.

Both negative cases follow by reversing the traversal: this exchanges left
and right turns and negates rotation by Lemma~\ref{lem:rot}(iii). Applying
the respective positive case to the reversed polygon proves each claim.
\end{proof}

\begin{remark}[the chirotope formula at a hull vertex]\label{prop:whitney}
For comparison with Whitney's formula~\cite{Whitney}, let $P$ be generic and suppose $\mu_1$ is a vertex of the convex hull of
$\{\mu_1,\ldots,\mu_n\}$. Then
\[
\rot(P)=\chi_{n,1,2}+\sum_{\substack{\{i,j\}\in X(P)\\ 1\le i<j\le n}}\chi_{i,i+1,j},
\]
where the crossing pairs are written with $i<j$ as integers in
$\{1,\ldots,n\}$ --- equivalently $E_i$ is visited before $E_j$ when the
traversal starts at $\mu_1$.
This is explanatory background only; no proof in this document uses this remark.
\end{remark}

\begin{remark}[the hull hypothesis cannot be dropped]\label{rem:whitney-base}
Without the hull hypothesis the formula can fail: for the embedded dart
$P=((2,1),(2,4),(0,0),(4,0))$, traversed counterclockwise, $\rot(P)=1$, there
are no crossings, and the turn at the reflex base vertex $\mu_1=(2,1)$ is
$-1$; the formula returns $-1$. With any of the three hull vertices as base it
returns $+1$. The
general formula with an arbitrary base vertex adds the winding numbers of the
two regions adjacent to the base vertex in place of the turn sign
\cite{Polyak}; at a hull vertex one adjacent region is the unbounded one,
with winding number $0$, and the general form reduces to the displayed one.
The hull hypothesis is sufficient, not necessary: at a convex vertex of an
embedded polygon the two adjacent winding numbers are $0$ and $\tau_1$
whether or not that vertex lies on the hull. The main text should either
restrict to a hull vertex or state the general form.
\end{remark}

\subsection{Cyclic shift and reversal}

\begin{lemma}[behaviour under shift and reversal]\label{lem:shift}
\begin{enumerate}
\item[(i)] The chirotope identities are
\[
 \chi_{ijk}(\sigma P)=\chi_{i+1,j+1,k+1}(P),\qquad
 \chi_{ijk}(\overline P)=\chi_{2-i,2-j,2-k}(P).
\]
The edge $E_i(\sigma P)$ is $E_{i+1}(P)$, and the edge
$E_i(\overline P)$ is $E_{1-i}(P)$ traversed backwards.
\item[(ii)] $\sigma$ and $\rho$ map $\mathcal U_n$ onto itself and map chambers
onto chambers; they preserve $X(P)$ up to relabelling of edges.
\item[(iii)] $\tau_i(\sigma P)=\tau_{i+1}(P)$ and $\tau_i(\overline P)=-\tau_{2-i}(P)$;
hence $\ell(\overline P)=n-\ell(P)-z(P)$, where $z(P)=\#\{i:\tau_i(P)=0\}$;
in particular $\ell(\overline P)=n-\ell(P)$ for $P\in\mathcal U_n$
(Lemma~\ref{lem:g1}(ii)).
\item[(iv)] for $P\in\mathcal R_n$: $\sigma P$ and $\overline P$ lie in
$\mathcal R_n$, $\rot(\sigma P)=\rot(P)$ and $\rot(\overline P)=-\rot(P)$.
\end{enumerate}
\status{proved (refereed: bench A 2026-09-12T06:44:52Z, bench B 2026-09-12T06:57:24Z; SM1-carried, read on SM12, repaired in round 12 (F-25-189, F-25-193), re-read on SM13)}
\end{lemma}
\begin{proof}
(i) is the definition; for the reversal, $\det(\mu_{2-j}-\mu_{2-i},\mu_{2-k}-\mu_{2-i})$
is the determinant of the triple $(2-i,2-j,2-k)$ of $P$. (ii): (G1) and (G2) are
conditions on the set of vertices and edges, which are preserved; both maps are
linear homeomorphisms of $(\RR^2)^n$. (iii): $\tau_i(\overline P)=
\chi_{i-1,i,i+1}(\overline P)=\chi_{3-i,2-i,1-i}(P)=-\chi_{1-i,2-i,3-i}(P)
=-\tau_{2-i}(P)$ by Lemma~\ref{lem:chi-basic}(i). Since $i\mapsto2-i$ is a
bijection of $\ZZ/n$, the turn multiset of $\overline P$ is the negative of
that of $P$: the left turns of $\overline P$ are the right turns of $P$, of
which there are $n-\ell(P)-z(P)$; on $\mathcal U_n$ every turn is $\pm1$
(Lemma~\ref{lem:g1}(ii)), so $z(P)=0$ there.
(iv): $\sigma$ permutes the edge vectors cyclically,
$\lt_i(\sigma P)=\lt_{i+1}(P)$, and $\rho$ negates and reverses them,
$\lt_i(\overline P)=-\lt_{1-i}(P)$, so the consecutive pairs
$(\lt_{i-1},\lt_i)$ of $\sigma P$ are those of $P$ and the consecutive pairs
of $\overline P$ are the pairs $(-\lt_k,-\lt_{k-1})$ of $P$; the condition of
Definition~\ref{def:regular} (every edge vector nonzero, no $\lt_k$ a negative
multiple of $\lt_{k-1}$) is symmetric in the pair and invariant under
negation, so it is preserved and $\sigma P,\overline P\in\mathcal R_n$. The
identity for $\rho$ is then Lemma~\ref{lem:rot}(iii), and for $\sigma$ the
multiset of principal turns is unchanged.
\end{proof}

\begin{definition}[admissible pairs]\label{def:admissible}
A pair $(n,r)\in\ZZ^2$ is \emph{admissible} if $n\geq3$, $2|r|<n$ and
$(n,r)\neq(3,0)$. It is \emph{minimal} if moreover $n=2|r|+1$ (for $r\neq0$)
or $(n,r)=(4,0)$. For an admissible pair the \emph{fibre} is
$\mathcal U_{n,r}=\{P\in\mathcal U_n:\rot(P)=r\}$; a vertex of a fibre is a
generic polygon, a point of the quotient.
\end{definition}

\begin{lemma}[nonempty fibres]
\label{lem:fibres}
For integers $n,r$, a generic polygon with $n$ vertices and rotation $r$
exists exactly when $n\geq3$, $2|r|<n$, and $(n,r)\ne(3,0)$.
Here genericity means that every determinant of three distinct vertices
is nonzero and no three distinct edge interiors meet. Moreover, for $n\geq3$
the generic polygons are dense in $(\RR^2)^n$: every nonempty open set of
labelled $n$-tuples contains a generic polygon.
\status{proved (refereed: bench A, 2026-09-05T23:45:25Z and 2026-09-06T01:58:43Z; transcribed from CV thm:zeroanchor(A) (C003); round~4 scope note: CV(A) supplies existence for CV's genericity (every relevant member of CV's guarded list nonzero, which constrains only vertex triples containing a cyclically consecutive pair), strictly weaker than the genericity stated here from $n=6$ on (bench~A's witness $((0,0),\allowbreak(8,-11),\allowbreak(2,0),\allowbreak(5,-9),\allowbreak(-7,0),\allowbreak(6,-11))$); the all-triple construction $D_{ijk}$ printed in the proof is this document's own strengthening, refereed sound by bench~A (F-25-138); round~5: the density clause its proof establishes (every nonempty open subset of $(\RR^2)^n$, $n\geq3$, contains a generic polygon) exported in the statement, with one sentence of the proof pointing the avoidance argument at an arbitrary open set, F-25-146)}
\end{lemma}
\begin{proof}
We use only Lemma~\ref{lem:rot}: rotation is
integer-valued and constant on paths in the regular locus; subdivision
preserves it; regular triangles have rotation $\pm1$; and a regular
$n$-gon satisfies $2|\operatorname{rot}|<n$. The regular locus consists
of polygons with nonzero edges and without antiparallel consecutive
directions. Thus the three stated conditions are necessary. We print
the existence argument, including its polynomial-avoidance step.

\emph{Regular seeds.}
First let $r>0$ and put $m=2r+1$, $\alpha=2\pi r/m$, and
$q=\exp(\mathrm i\alpha)$. Identify the plane with $\mathbb C$ and take
the ordered vertices $z_j=q^{j-1}$ for $1\leq j\leq m$.
The integers $r,m$ are coprime, since a common divisor would divide
$m-2r=1$. These vertices are therefore distinct. With cyclic indices,
the edge directions satisfy
\begin{equation}\label{fibpr:directions}
d_j=z_{j+1}-z_j=q^{j-1}(q-1).
\end{equation}
Here $q^m=1$, so this includes the closing edge, and $q\ne1$, so every
edge is nonzero. Successive directions satisfy $d_j=q d_{j-1}$, also
across the closing corner. Since $0<\alpha<\pi$, every principal turn
is exactly $\alpha$. The seed is regular, and its total turn is
\begin{equation}\label{fibpr:rotation}
\sum_{j=1}^{m}\vartheta_j=m\alpha=2\pi r.
\end{equation}
Consequently its rotation is $r$. For negative $r$, reverse the seed
constructed for $|r|$; reversal negates rotation and preserves regularity.
This construction uses no genericity or concurrency claim about a
regular star and does not refer to a later star definition.

For $r=0$, use the four ordered vertices
$(0,0),(2,2),(0,2),(2,0)$. Its four directions are
$(2,2),(-2,0),(2,-2),(-2,0)$, and its principal turns, starting at the
first vertex, are respectively
\begin{equation}\label{fibpr:bowtie}
-3\pi/4,\quad 3\pi/4,\quad 3\pi/4,\quad-3\pi/4.
\end{equation}
Their sum is zero and none is antiparallel. This is a regular seed of
rotation zero. In each case let $m$ be the number of seed vertices.
The stated admissibility conditions imply $n\geq m$.

\emph{Subdivision and a regular neighbourhood.}
Insert $n-m$ distinct interior subdivision points in traversal order
on any seed edge. All refined edge vectors are positive multiples of
the former vector; each new principal turn is zero and all old turns
are unchanged. The resulting $n$-gon $P_0$ is regular with rotation $r$.
The regular locus is open: nonzero vectors form an open set, and for
two nonzero vectors the condition that their unit directions sum to a
nonzero vector is open and excludes exactly antiparallel directions.
Choose an open ball $B$ about $P_0$ contained in this locus.
Every point of $B$ is connected to $P_0$ by a straight path in $B$,
so it has rotation $r$ by the preceding rotation-number lemma.

\emph{Explicit nonzero polynomials.}
All coordinates of all $n$ vertices may vary in $B$.
For every three distinct indices $i,j,k$, let
\begin{equation}\label{fibpr:triple}
D_{ijk}=\det(\mu_j-\mu_i,\mu_k-\mu_i).
\end{equation}
This polynomial is not identically zero: assign these three vertices
the coordinates $(0,0),(1,0),(0,1)$, respectively, obtaining value $1$.
The remaining coordinates can be arbitrary for this test.

For an edge from $(x_i,y_i)$ to $(x_{i+1},y_{i+1})$, define its
homogeneous line-coefficient row by
\begin{equation}\label{fibpr:line}
\ell_i=(y_i-y_{i+1},\ x_{i+1}-x_i,\
               x_i y_{i+1}-x_{i+1}y_i).
\end{equation}
Both endpoints satisfy $\ell_i\cdot(x,y,1)=0$ by direct substitution.
For each unordered triple of pairwise remote edges $i,j,k$, put
\begin{equation}\label{fibpr:concurrence}
H_{ijk}=\det\begin{pmatrix}\ell_i\\\ell_j\\\ell_k\end{pmatrix}.
\end{equation}
These three edges have six distinct endpoint indices. Hence we can
assign their ordered endpoint pairs independently as
$(0,0),(1,0)$; $(0,0),(0,1)$; and $(0,1),(1,1)$.
The rows in this assignment are $(0,1,0)$, $(-1,0,0)$, and $(0,1,-1)$;
their determinant is $-1$. Thus each $H_{ijk}$ is a nonzero polynomial.
Coinciding coordinates in this witness are harmless: it proves a
polynomial identity is nonzero on the full coordinate space, rather
than purporting to be a generic polygon or a point of $B$.

If the three supporting lines have a common finite point $(x,y)$,
their coefficient matrix annihilates the nonzero vector $(x,y,1)$.
Its determinant then vanishes. Thus $H_{ijk}\ne0$ excludes triple
interior concurrence for these three edges. It is a sufficient
condition; no converse is used.

\emph{Simultaneous avoidance in the small ball.}
Let $F$ be the finite product of all the polynomials $D_{ijk}$, taking
one order for each index triple, and all the polynomials $H_{ijk}$.
An empty family of edge triples contributes the factor $1$.
The polynomial ring over $\mathbb R$ has no zero divisors (the highest
nonzero monomials in lexicographic order multiply to a nonzero highest
monomial), so $F$ is nonzero.
For completeness, a nonzero real polynomial cannot vanish on a
nonempty open box. For one variable this follows from the finite root
bound. Induct on the number of variables, writing the polynomial as a
polynomial in the last variable with coefficient polynomials in the
others. If it vanishes on a box, fixing the other variables makes it
vanish on an interval; all its coefficients vanish at that fixed
choice. They therefore vanish on the lower-dimensional box, and the
inductive hypothesis makes every coefficient identically zero, a
contradiction. Every open ball contains an open box. Hence some
$P\in B$ has $F(P)\ne0$. The same avoidance applies to any nonempty open
subset of $(\RR^2)^n$ in place of $B$, since it too contains an open box;
with the genericity of such a $P$ shown next, this proves the density clause.

At this $P$, the nonvanishing of every $D_{ijk}$ gives (G1), in
particular nonzero edges and noncollinear consecutive vertices.
Adjacent edges can meet only at their common endpoint: their lines
are distinct and nonparallel because the consecutive-vertex
determinant is nonzero. A point in three edge interiors therefore
would require three pairwise remote edges. Its concurrence would
force the corresponding $H_{ijk}$ to vanish, contrary to $F(P)\ne0$.
This proves (G2). The polygon is generic and has rotation $r$ because
it lies in $B$. This proves sufficiency and the lemma.
\end{proof}

\subsection{Walls}

\begin{definition}[wall germ; sides]\label{def:germ}
A \emph{wall germ} is a continuous map $P\colon(-\varepsilon,\varepsilon)\to
(\RR^2)^n$ such that $P(t)\in\mathcal U_n$ for $t\neq0$ and $P(0)\notin
\mathcal U_n$. Its \emph{sides} $P_-$ and $P_+$ are the chambers containing
$P((-\varepsilon,0))$ and $P((0,\varepsilon))$; each of these two sets is
connected and generic, hence lies in one chamber. For a function $F$ that is
constant on chambers, $F(P_\pm)$ denotes its value on the respective side.
At the centre $P(0)$ put
\begin{align*}
Z_{\rm pt}&=\bigl\{\{i,j,k\}\text{ pairwise distinct}:\chi_{ijk}(P(0))=0\bigr\},\\
Z_{\rm c}&=\bigl\{\{e,f,g\}\text{ pairwise remote}:
\operatorname{relint}E_e\cap\operatorname{relint}E_f\cap\operatorname{relint}E_g\neq\varnothing\text{ at }t=0\bigr\}.
\end{align*}
A real function $\phi$ of the germ \emph{changes sign at $0$} if there is
$\delta>0$ with $\phi(P(t))\phi(P(-t))<0$ for $0<t<\delta$.
\end{definition}

\begin{lemma}[sides of a triple wall]\label{lem:triple-sides}
Let a wall germ have $Z_{\rm pt}=\varnothing$ and
$Z_{\rm c}=\{\{e,f,g\}\}$. Then for small $t\neq0$ the three pairs
$\{e,f\},\{e,g\},\{f,g\}$ are crossings, their three crossing points
are distinct, and on each of $E_e,E_f,E_g$ its two crossings occupy adjacent
positions among the crossing visits on that edge.
\status{proved (refereed: bench A, 2026-09-05T15:51:04Z; transcribed from RA R-LOC-2, stated there for a simple transversal event; the printed proof covers every germ with $Z_{\rm pt}=\varnothing$ and $Z_{\rm c}$ the single triple $\{e,f,g\}$ (C015; F-25-44))}
\end{lemma}
\begin{proof}
At the centre (G1) holds. In particular all vertices and edges are nonzero,
adjacent edges meet only at their common endpoint, and no vertex lies on the
line of a nonincident edge. The specified common point $q$ lies in all three
relative interiors. Two of their directions cannot be parallel: their lines
would then coincide, forcing an endpoint of one remote edge onto the line
of the other, contrary to (G1). Thus all three central intersections are
transverse. Lemma~\ref{lem:wall-segment-stability} makes them persistent
interior crossings. On either punctured side (G2) makes the three points
distinct.

In fact every remote edge pair at the centre is either disjoint or a transverse
interior pair: endpoint incidence or overlapping collinear remote edges would
contradict (G1). Thus the same stability lemma applies to all remote pairs and
gives a fixed finite crossing set with continuous parameters near the centre.

Suppose, contrary to the required adjacency on $E_e$, that for a sequence
$t_\nu\to0$ another crossing lies between $x_{ef}$ and $x_{eg}$ on $E_e$.
There are finitely many crossing labels, so after taking a subsequence this
crossing has one fixed other edge $h\notin\{e,f,g\}$. The two bounding
parameters converge to the parameter of $q$; the intervening parameter must
have the same limit. Continuity of its intersection parameters puts $q$ in
the relative interiors of $E_e$ and $E_h$. It is also interior to $E_f$.
The edges $E_h,E_f$ cannot be adjacent: under (G1) adjacent edges meet only
at their shared endpoint, whereas $q$ is interior to both. Hence $e,f,h$
are pairwise remote and define an element of $Z_{\rm c}$ distinct from
$\{e,f,g\}$, a contradiction. The same argument on $E_f$ and $E_g$ proves
all three adjacency assertions on one common punctured interval. No sign-change
hypothesis was used for this lemma.
\end{proof}


\begin{definition}[the named walls]\label{def:walls}
A wall germ is \emph{simple of one of the following types}; the types are
mutually exclusive.
\begin{enumerate}
\item[(F)] \emph{Flat at $j$} ($n\geq4$): $Z_{\rm pt}=\{\{j-1,j,j+1\}\}$,
$Z_{\rm c}=\varnothing$, $\mu_j(0)$ lies strictly between $\mu_{j-1}(0)$ and
$\mu_{j+1}(0)$, and $\tau_j$ changes sign at $0$. The \emph{right side} is the
side with $\tau_j=-1$ and the \emph{left side} the one with $\tau_j=+1$.
\item[(K)] \emph{Cusp at $j$} ($n\geq4$): $Z_{\rm pt}=\{\{j-1,j,j+1\}\}$,
$Z_{\rm c}=\varnothing$, the three points $\mu_{j-1}(0),\mu_j(0),\mu_{j+1}(0)$
are collinear with $\mu_j(0)$ outside the closed segment
$[\mu_{j-1}(0),\mu_{j+1}(0)]$, and $\tau_j$ changes sign at $0$. The
\emph{newborn pair} is $\{j-1,j+1\}$ if $\mu_{j+1}(0)$ lies between
$\mu_{j-1}(0)$ and $\mu_j(0)$, and $\{j-2,j\}$ if $\mu_{j-1}(0)$ lies between
$\mu_j(0)$ and $\mu_{j+1}(0)$. The \emph{loop side} $P_{\rm loop}$ is the side
on which the newborn pair is a crossing and the \emph{no-loop side}
$P_{\rm no}$ the other. The cusp is
\emph{empty} if on the loop side the two visits of the newborn crossing are
cyclically adjacent in the Gauss word.
\item[(V)] \emph{Vertex--edge at $(M;a)$}: $M\notin\{a-1,a,a+1,a+2\}$,
$Z_{\rm pt}=\{\{a,a+1,M\}\}$, $Z_{\rm c}=\varnothing$, $\mu_M(0)$ lies in the
relative interior of $E_a(0)$, and $\chi_{a,a+1,M}$ changes sign at $0$. The
germ is of \emph{bigon} type if $\chi_{a,a+1,M-1}(P(0))=\chi_{a,a+1,M+1}(P(0))$
and of \emph{sliding} type otherwise. The \emph{contact sign} is
$s=\chi_{a,a+1,M}(P_-)$.
\item[(T)] \emph{Triple at $\{e,f,g\}$}: $Z_{\rm pt}=\varnothing$,
$Z_{\rm c}=\{\{e,f,g\}\}$, and on each of the three edges the two crossing
parameters of the other two edges (which exist on both sides by
Lemma~\ref{lem:triple-sides}) have a difference that changes sign at $0$.
\item[(E)] \emph{Exterior extension at $(M;a)$}: $M\notin\{a-1,a,a+1,a+2\}$,
$Z_{\rm pt}=\{\{a,a+1,M\}\}$, $Z_{\rm c}=\varnothing$, $\mu_M(0)$ lies on the
line of $E_a(0)$ outside the closed segment $E_a(0)$, and $\chi_{a,a+1,M}$
changes sign at $0$.
\item[(C)] \emph{Pure cut at $\{i,j,k\}$}: no two of $i,j,k$ are consecutive
modulo $n$, $Z_{\rm pt}=\{\{i,j,k\}\}$, $Z_{\rm c}=\varnothing$, and
$\chi_{ijk}$ changes sign at $0$.
\end{enumerate}
The walls (E) and (C) are called \emph{silent}. The type is determined by the
centre; the sign-change conditions exclude tangential germs.
\end{definition}

\begin{lemma}[flat-side geometry]\label{lem:flat-sides}
Let $n\geq4$ and let $P(t)=(\mu_1(t),\ldots,\mu_n(t))$ be continuous for
$|t|<\varepsilon$. Indices are cyclic modulo $n$; put
$E_i=[\mu_i,\mu_{i+1}]$ and $d_i=\mu_{i+1}-\mu_i$.
For $t\ne0$ assume that no three distinct vertices are collinear and no
three edge interiors concur. At $t=0$ suppose that the only collinear
triple is $\{j-1,j,j+1\}$, that no three pairwise remote edge interiors
concur, and that $\mu_j$ is strictly between $\mu_{j-1}$ and
$\mu_{j+1}$. Suppose also that
$\tau_j(t)=\operatorname{sgn}\det(d_{j-1}(t),d_j(t))$ changes sign at zero.
At a transverse crossing of $E_e,E_f$, the \emph{positive over/under
data} mean that $E_e$ is over $E_f$ exactly when $\det(d_e,d_f)>0$.
Then, after shrinking the interval:
\begin{enumerate}
\item[(i)] The central polygon has nonzero edges and no antiparallel
consecutive directions. Its only zero turn is at $j$, where the two
directions are positive multiples. Every other point-triple determinant
has constant nonzero sign on the whole interval.
\item[(ii)] The set of remote edge pairs that cross is constant,
including at zero. Each crossing is interior and transverse, no crossing
is a vertex, all crossing points are distinct, and the crossing visits
have constant order on every edge. Hence their cyclic Gauss word,
pairing, interlacement graph, and positive over/under data are constant.
Here the central crossing data are read from the segments of $P(0)$
exactly as for a generic polygon: a crossing is a transverse intersection
of the relative interiors of two remote edge segments, its two visits are
ordered along each edge by the intersection parameters, and the Gauss
word, pairing, interlacement graph and positive over/under data are formed
from these visits as in Definitions~\ref{def:crossings}, \ref{def:gauss}
and~\ref{def:interlace}; no function on the generic locus is evaluated at
the flat centre.
\item[(iii)] The deletion $Q=P(0)\setminus j$ is generic. The directions
$u=d_{j-1}(0)$ and $v=d_j(0)$ fuse to $w=\mu_{j+1}(0)-\mu_{j-1}(0)$,
and there is $0<\lambda<1$ with
\begin{equation}\label{flatpr:fusion}
u=\lambda w,\qquad v=(1-\lambda)w.
\end{equation}
Replacing either incident edge in a crossing label by the fused edge
gives a bijection with the crossings of $Q$, preserving their cyclic
visit order, pairing, determinant signs, and positive over/under data.
The deletions $P(t)\setminus j$ are generic on a smaller interval and
lie in the chamber of $Q$.
\end{enumerate}
\status{proved (refereed: bench A, 2026-09-05T23:45:25Z and 2026-09-06T04:22:56Z; transcribed from CV lem:flatdata --- (i),(ii) from its clauses (i),(ii); (iii) from its clause (ii) (the deletion is generic and carries the same data under the fused-edge identification) together with the fused-segment paragraph of its proof (the incident directions are positive multiples of the fused direction) and, for the nearby deletions lying in the chamber of the central one, CV thm:main(A)(i) (d9\_induction; stated there for CV's simple flat events under its guarded genericity, whereas this lemma asserts it under the all-triple genericity of its hypothesis, the printed proof carrying the wider scope; bench B, P-25-37/38, F-25-177) (P-25-3, C007); round~4: the flat-centre crossing data defined in clause~(ii), F-25-140; the proof's ``positive lift'' phrase replaced by the statement's positive over/under data, F-25-112)}
\end{lemma}

\begin{proof}
Write $A=\mu_{j-1}(0)$, $M=\mu_j(0)$ and $B=\mu_{j+1}(0)$.
All central vertices are distinct: if two coincided, adjoining each
of two further indices would produce two distinct collinear triples;
there are two further indices because $n\geq4$. This contradicts the
unique-triple assumption. Thus all edges are nonzero. Every turn other
than the one at $j$ uses a different triple, so its determinant is
nonzero. Strict betweenness gives $M=A+\lambda(B-A)$ for exactly one
$\lambda\in(0,1)$, proving \eqref{flatpr:fusion}. In particular the
turn at $j$ is positive-flat, not antiparallel. Continuity and the
finitely many nonzero central determinants give (i).

\emph{No vertex contact at the centre.}
A vertex on the line of an edge not incident to it would give a
collinear triple. The sole permitted triple can occur in this way only
as $B$ on the line of $[A,M]$ or $A$ on the line of $[M,B]$.
Strict betweenness puts each such point outside the corresponding
closed segment. Thus no vertex lies on a nonincident closed segment.
Two adjacent edges other than the flat pair have nonparallel directions,
and meet only at their common endpoint. The flat pair consists of the
successive subsegments $[A,M]$ and $[M,B]$, also meeting only at their
common endpoint. Hence no adjacent edge interiors meet.

\emph{Crossing pairs.}
For a remote pair whose four endpoint-against-line triples avoid
$\{j-1,j,j+1\}$, those determinants have nonzero limits and constant
signs. Two segments with all endpoints off the other segment's line
cross if and only if each line separates the other segment's two
endpoints. This is the pair of strict opposite-sign tests. To see
sufficiency, intersect the two lines: opposite endpoint signs place
their unique intersection inside each segment. Necessity follows by
restricting the affine signed line equation to the other segment.
The tests therefore give the same crossing answer at zero and on both
sides. An actual intersection cannot be parallel: parallel intersecting
lines would coincide, forcing an endpoint determinant to vanish.

The only remote pairs whose tests can use the critical triple are
$\{j-2,j\}$ and $\{j-1,j+1\}$. The line of $E_{j-2}(0)$ meets the
flat line only at $A$; its other endpoint is off that line by the
unique-triple assumption. Since $A\notin[M,B]$, the compact segments
$E_{j-2}(0)$ and $E_j(0)$ are disjoint. Similarly the line of
$E_{j+1}(0)$ meets the flat line only at $B\notin[A,M]$, so
$E_{j-1}(0)$ and $E_{j+1}(0)$ are disjoint. Positive distances between
these compact segment pairs persist under small changes of their
endpoints. These pairs remain noncrossing on both sides. This covers
all remote pairs, including when $n=4$.

\emph{Order and crossing records.}
Each crossing at zero has two nonparallel directions. Solving its
two-by-two line-intersection system therefore gives continuous crossing
parameters, strictly between zero and one at zero. There are finitely
many pairs, so all these strict inequalities persist on one interval.
If two distinct central crossings had the same point, at least three
edge interiors would contain that point. None of those edges could be
adjacent, by the preceding paragraph. This would be a forbidden
concurrence of three pairwise remote edges. Thus crossing points are
distinct. In particular, on a fixed edge all crossing parameters are
distinct. The finitely many nonzero parameter differences retain their
signs by continuity; crossing visits neither collide nor reorder.
The Gauss word is just these edgewise lists concatenated in traversal
order. Pairing is indexed by the fixed edge pairs, and interlacement is
read from that word. For each actual crossing, the nonzero direction
determinant keeps its sign. The same physical strand therefore remains
over in the positive over/under data of the statement. This proves (ii)
without a triple-wall lemma.

\emph{Deletion and fusion.}
The deletion's point triples are central parent triples avoiding $j$,
so they are all noncollinear. At zero the oriented segment $[A,B]$
is precisely the traversal of $[A,M]$ followed by $[M,B]$.
No crossing of a remote edge with this segment can occur at $M$:
that would put the parent vertex $M$ on a nonincident edge, already
excluded. Thus such a crossing lies in exactly one of the two open
subsegments and inherits its transverse direction and its pairing.
An edge remote from $[A,B]$ is remote from both old subsegments.
The reverse implication needs care only for the neighbours at $A,B$:
their possible opposite-subsegment pairs are exactly the two compact
disjoint pairs treated above. They create no discarded crossing.
One remote straight edge cannot meet both old subsegments, since it
meets their common line at most once. Hence the crossing-label map is
bijective.

If a crossing parameter on $[A,M]$ is $a$, its parameter on $[A,B]$
is $\lambda a$; if it is $b$ on $[M,B]$, its fused parameter is
$\lambda+(1-\lambda)b$. Thus the old first-edge visits precede the
old second-edge visits, in exactly their original orders. Equation
\eqref{flatpr:fusion} shows that every direction determinant involved
changes only by a positive factor, preserving signs and over/under
data. No triple of deletion edge interiors can concur: away from $M$
it would give three parent edge interiors at the same point; at $M$
it would put a remote parent edge through that parent vertex. Both
are excluded. This proves genericity and all stated crossing data for
$Q$. Finally, the finitely many strict noncollinearity, transverse
crossing, disjointness and crossing-order conditions used above also
prove openness of genericity near $Q$. By continuity, the connected
small path $P(t)\setminus j$ stays in that open generic neighbourhood
and hence in the chamber of $Q$. This proves (iii).
\end{proof}

\begin{lemma}[the sides of the named walls]\label{lem:wall-sides}
Let the germ be simple of type (F), (V), (T), (E) or (C). Then
$P(0)\in\mathcal R_n$. After shrinking the interval, every chirotope outside
$Z_{\rm pt}$ is nonzero with constant sign. Every crossing not involved in
the stated local degeneration persists, and the order of any two persistent
visits with distinct central parameters is constant. Moreover:
\begin{enumerate}
\item[(F)] The crossing set and Gauss word are the same on both sides and
at the centre.
\item[(V)] In the bigon branch,
\[
 X(P_+)\mathbin\triangle X(P_-)=\{\{a,M-1\},\{a,M\}\},
\]
with both pairs crossing on one side and neither on the other. In the sliding
branch exactly one of these two pairs crosses on each side, and the two sides
exchange them. No other crossing is created or lost, and no order changes
outside the contact neighbourhood.
\item[(T)] The crossing set is the same on both sides. The only visit-order
changes are the exchanges of the two triangle crossings on each of
$E_e,E_f,E_g$.
\item[(E),(C)] The crossing set and Gauss word are the same on both sides
and at the centre.
\end{enumerate}
Here, when crossing data are mentioned at a nongeneric centre, they mean the
transverse intersections of the relative interiors of the actual segments and
their traversal order. This does not evaluate a function defined only on
generic polygons at that centre.
\status{proved (refereed: bench A, 2026-09-05T15:51:04Z; transcribed from CV lem:flatdata; CV lem:silence, its proof, for (E),(C); RC tr:five-events, gw:germs (S3/S4/S7/R3 only), flat:pure (C015; F-25-45))}
\end{lemma}
\begin{proof}
Case (F), including its central regularity and crossing data, is
Lemma~\ref{lem:flat-sides}. We prove the remaining cases directly
in the one-based tail convention $d_i=\mu_{i+1}-\mu_i$.

\emph{Central distinctness and regularity.}
The remoteness requirements give $n\geq5$ in (V),(E) and $n\geq6$ in (T),(C).
If two central vertices coincided, adjoining each of two further indices
would give two different zero point triples. This contradicts the permitted
zero set, which has at most one member. Thus all central vertices are distinct
and every edge is nonzero. A zero turn would give a consecutive zero triple.
There is no such triple in any of (V),(T),(E),(C): the point zero set is empty
in (T), has no adjacent pair in (C), and has just the adjacent pair $a,a+1$
in (V),(E), because $M\notin\{a-1,a,a+1,a+2\}$. Hence every central turn
in these cases is nonzero. In particular the centre has no kink and belongs
to $\mathcal R_n$. Each other point determinant is nonzero at zero and is
continuous; finitely many such determinants retain their signs nearby.

\emph{Which segment pairs can change.}
No remote pair can have collinear overlapping segments at the centre: its
four distinct endpoint indices on a line would give more than one zero
point triple. Apart from an endpoint incidence, a remote pair is therefore
either disjoint or a transverse interior pair, and
Lemma~\ref{lem:wall-segment-stability} applies. An endpoint incidence requires
a zero point triple containing the two endpoints of its base edge. In (C)
the only zero triple contains no adjacent pair, and in (T) there is no zero
triple, so these cases have no endpoint incidence. In (V),(E), the only
possible incidence is the specified vertex $M$ on the line of $E_a$.
The only possibly affected remote pairs are consequently
$\{a,M-1\}$ and $\{a,M\}$. All other pairs satisfy the stability lemma.

\emph{The two vertex--edge branches, with the finite-segment test.}
Write $A=\mu_a$, $B=\mu_{a+1}$, $D=B-A$, $M=\mu_M$, and take in turn
$U=\mu_{M-1}$ and $U=\mu_{M+1}$. Put
\begin{equation}\label{wallgeo:heights}
 h=\det(D,M-A),\qquad h_U=\det(D,U-A).
\end{equation}
At the centre $h=0$ and $h_U\ne0$: another zero $h_U$ would give a second
point triple, since both neighbour indices differ from $a,a+1,M$.
Shrink so $h_U$ and $h_U-h$ are nonzero. The intersection of the supporting
line of $E_a$ with the line $MU$ has the exact form
\begin{equation}\label{wallgeo:contact-foot}
 q_U=M+r_U(U-M),\qquad r_U=-\frac{h}{h_U-h}.
\end{equation}
Indeed its height is $h+r_U(h_U-h)=0$. When $h$ and $h_U$ have opposite
signs, $r_U=|h|/(|h|+|h_U|)$ lies strictly between zero and one. When their
signs agree, the height of every point in the relative interior of $MU$ has
that common sign, so there is no such intersection. In the first case
$r_U\to0$ and $q_U\to M(0)$. In case (V), $M(0)$ is strictly inside $E_a(0)$.
The parameter of $q_U$ on the line $AB$ therefore tends to a number in $(0,1)$
and remains in $(0,1)$ nearby. This supplies the second, finite-segment
condition, not merely the straddling of the supporting line. The crossing is
transverse since the two line directions have determinant $h_U-h\ne0$.

Consequently each contact pair crosses precisely when the sign of $h$ is
opposite to its neighbour sign. The two neighbour signs are equal in the
bigon branch, giving both crossings or neither. They are opposite in the
sliding branch, giving precisely one crossing on each side. By the defining
sign change of $h$, the asserted exchanges follow. The predecessor leg is
$E_{M-1}$ and the successor leg is $E_M$, so these are exactly the displayed
tail-indexed pairs.

\emph{Exterior extension.}
At the centre in (E), each contact leg meets the supporting line of $E_a$
only at $M(0)$, because $h_U(0)\ne0$. That point is outside the closed
segment $E_a(0)$. Thus each contact leg is compact-disjoint from $E_a(0)$,
and both pairs remain disjoint by Lemma~\ref{lem:wall-segment-stability}.
All remote crossing pairs are therefore constant, including at the centre,
in (E), as they already are in (C) and (T).

\emph{Orders and localization.}
Consider two persistent crossings on a common edge whose central parameters
coincide. The other two edges then meet that base edge at the same point.
If those other edges are remote, their common interior point gives a member
of $Z_{\rm c}$. If they are adjacent, their nonzero central turn implies that
their lines meet only at their shared vertex. The coincidence would put that
vertex on the base edge. In (E),(C) no such endpoint incidence exists; in
(T) none exists either, and its only concurrence is the named triple.
It follows that all parameters in (E),(C) are centrally distinct and their
orders are constant. Their fixed edge order then gives a fixed cyclic Gauss
word, including at zero. In (T) the only possible ties are the two triangle
visits on a bundle edge. Lemma~\ref{lem:triple-sides} makes each such pair
adjacent. Their differences change sign by the definition of (T), so they
exchange order, and every other order is fixed by the stability lemma.

For (V), no persistent crossing has central point $M(0)$: any additional
edge through that point would give a second endpoint-incidence triple.
All persistent crossing parameters on $E_a$ are therefore separated from its
contact parameter. Those on either incident leg are separated from its
endpoint at $M$. Formula~\eqref{wallgeo:contact-foot} puts any newborn contact
crossings arbitrarily close to those contact parameters. Finiteness gives
neighbourhoods free of other crossings. Every other pair of persistent visits
has distinct central parameters, since a tie would give either a forbidden
triple concurrence or an additional vertex incidence. Their orders remain
fixed. This proves the claimed localization and completes all cases.
\end{proof}



\begin{lemma}[sides of a cusp]\label{lem:cusp-sides}
Let $n\geq4$ and let
$P(t)=(\mu_1(t),\ldots,\mu_n(t))$, $|t|<\varepsilon$, be continuous.
Indices are read modulo $n$, and
$E_i=[\mu_i,\mu_{i+1}]$,
$\widetilde\lambda_i=\mu_{i+1}-\mu_i$ and
$\tau_i=\operatorname{sgn}\det(\widetilde\lambda_{i-1},\widetilde\lambda_i)$.
Assume that $P(t)$ is generic for $t\ne0$: no three distinct vertices are
collinear and no three edge interiors concur. At $t=0$ assume that the only
collinear triple of distinct vertex indices is $\{j-1,j,j+1\}$, that no
three pairwise remote edge interiors concur, that $\mu_j$ lies outside the
closed segment $[\mu_{j-1},\mu_{j+1}]$, and that $\tau_j$ changes sign.
These are the simple-cusp hypotheses in the present conventions.

Use the following data, with the betweenness conditions imposed at $t=0$:
\begin{equation}\label{cusppr:data}
\begin{array}{c|c|c|c|c}
\text{case}&\text{middle point on the line}&(f,g)&(c_1,c_2)&\Delta(t)\\ \hline
\mathrm A&\mu_{j+1}\in(\mu_{j-1},\mu_j)&(j-1,j+1)&(j,j+1)&
 \det(\widetilde\lambda_{j-1},\widetilde\lambda_{j+1})\\
\mathrm B&\mu_{j-1}\in(\mu_j,\mu_{j+1})&(j-2,j)&(j-1,j)&
 \det(\widetilde\lambda_{j-2},\widetilde\lambda_j)
\end{array}
\end{equation}
Exactly one case applies. After shrinking the parameter interval:
\begin{enumerate}
\item[(i)] $\Delta$ is nonzero with constant sign. The remote pair
$\{f,g\}$ crosses precisely on the side with
$\tau_j=-\operatorname{sgn}\Delta$, called the loop side; it does not cross
on the side with $\tau_j=\operatorname{sgn}\Delta$, called the no-loop side.
No other crossing pair changes.
\item[(ii)] On the no-loop side the turns at $c_1,c_2$ have opposite signs.
On the loop side both have sign $\tau_j$.
\item[(iii)] On the loop side, one oriented traversal arc between the two
visits of the newborn crossing consists of a terminal portion of $E_f$,
the complete edge between $c_1$ and $c_2$, and an initial portion of $E_g$.
Its only polygon vertices are $\mu_{c_1},\mu_{c_2}$. This arc is allowed to
contain other crossing visits. If the two newborn visits are cyclically
adjacent in the Gauss word, this particular arc contains no other crossing
visit. In that case the complete edge between $c_1,c_2$ carries no crossing
on either side of the wall.
\item[(iv)] With rotation defined by the sum of principal turns,
\begin{equation}\label{cusppr:rotation-law}
\operatorname{rot}(P_{\rm loop})-\operatorname{rot}(P_{\rm no})
=\tau_j(P_{\rm loop})\in\{-1,1\}.
\end{equation}
\end{enumerate}
\status{proved (refereed: bench A, 2026-09-05T15:51:04Z; transcribed from CV lem:cuspcases; RC mt:s4 for the rotation clause, $\kappa=+\tau_j$ on the loop side (P-25-1, C004))}
\end{lemma}

\begin{proof}
First, all vertex coordinates at the centre are distinct. If two coincided,
their triples with either of two further vertex indices would both be
collinear. Such further indices exist because $n\geq4$, contradicting the
sole-triple hypothesis. All edges at the centre are therefore nonzero.
Put $A=\mu_{j-1}$, $M=\mu_j$, $B=\mu_{j+1}$, and, as functions of $t$, put
\begin{equation}\label{cusppr:local-data}
u=M-A,\qquad v=B-M,\qquad T=\det(u,v).
\end{equation}
Thus $\tau_j=\operatorname{sgn}T$ on the punctured sides. The distinct
collinear points $A,M,B$ have a unique middle point; it is not $M$.
This proves the two-case assertion. On each punctured side every point-triple
determinant is nonzero, so its sign is constant by continuity.

\emph{The newborn crossing in case A.}
Let $w=\widetilde\lambda_{j+1}$, so that $E_g$ starts at $B$ and has
direction $w$. Here $\Delta=\det(u,w)$. If $\Delta(0)=0$, the nonzero
vector $w(0)$ would be parallel to the cusp line through $A,M,B$.
Its other endpoint $\mu_{j+2}(0)$ would lie on that line too. The indices
$j-1,j,j+2$ are distinct and give a second zero triple, a contradiction.
Therefore $\Delta(0)\ne0$; it has one nonzero sign on a smaller interval.

Write the intersection of the two supporting lines as $X=B+s w$.
Membership in the line through $E_f$ means $\det(u,X-A)=0$.
Since $B-A=u+v$, bilinearity gives, in separate steps,
\begin{equation}\label{cusppr:A-det}
\det(u,B-A)=\det(u,u)+\det(u,v)=T,
\end{equation}
\begin{equation}\label{cusppr:A-intersection}
0=\det(u,X-A)=T+s\Delta,
\qquad s=-\frac{T}{\Delta}.
\end{equation}
At $t=0$, $X=B$ lies strictly inside $E_f=[A,M]$ and $s=0$.
The line intersection varies continuously because $\Delta$ remains nonzero.
Consequently its parameter on $E_f$ remains strictly between $0$ and $1$,
and $|s|<1$ on a smaller interval. It lies in the interior of $E_g$ exactly
when $s>0$. Thus $E_f,E_g$ cross exactly when $T$ and $\Delta$ have
opposite signs. This also proves that the new crossing is interior and
transverse on the stated side.

\emph{The newborn crossing in case B.}
Let $h=\widetilde\lambda_{j-2}$, so that $E_f$ ends at $A$ and has
direction $h$. Here $\Delta=\det(h,v)$. If $\Delta(0)=0$, the other
endpoint $\mu_{j-2}(0)$ of $E_f$ would also lie on the cusp line, again
giving a second zero triple. Hence $\Delta$ is nonzero with constant sign
on a smaller interval. Write the intersection as $X=A-s h$, so $s>0$
means moving into $E_f$ from its endpoint $A$. The relevant determinant is
\begin{equation}\label{cusppr:B-det}
\det(v,A-M)=\det(v,-u)=-\det(v,u)=T.
\end{equation}
Because $\det(v,h)=-\Delta$, its line equation becomes
\begin{equation}\label{cusppr:B-intersection}
0=\det(v,X-M)=T-s\det(v,h)=T+s\Delta,
\qquad s=-\frac{T}{\Delta}.
\end{equation}
At $t=0$, $X=A$ is strictly inside $E_g=[M,B]$. Its parameter there
remains interior, and $|s|<1$ on a smaller interval. The two segments
therefore cross exactly when $s>0$, equivalently when $T$ and $\Delta$
have opposite signs. Since $T$ changes sign, the loop and no-loop sides
in both cases are distinct and have exactly the signs in (i).

\emph{No other crossing pair changes.}
For any remote pair, strict crossing is determined by the four signs of the
endpoint-against-line determinants: each segment must separate the endpoints
of the other. This criterion applies on both generic punctured sides.
If none of those four triples is $\{j-1,j,j+1\}$, all four signs have
the same nonzero limits and hence agree across the wall. Such a crossing
pair cannot change.

The only remote pairs whose tests can contain that triple are
$\{j-1,j+1\}$ and $\{j-2,j\}$. Indeed the only edges joining two critical
vertices are $[A,M]$ and $[M,B]$; adjoining the third critical vertex as an
endpoint of a remote edge gives precisely these two pairs. In case A,
the unused pair is $E_{j-2},E_j$. At the centre the supporting line of
$E_{j-2}$ meets the cusp line only at $A$, since its other endpoint is
off that line. But $A$ lies outside $[M,B]$, so the two compact segments
are disjoint and remain disjoint nearby. In case B the unused pair is
$E_{j-1},E_{j+1}$; the supporting line of $E_{j+1}$ meets the cusp line
only at $B$, which lies outside $[A,M]$. The same compact-disjointness
argument applies. This proves the last assertion of (i). It does not
assert that the order of thread crossings is unchanged.

\emph{The two needle turns.}
In case A there is $a>0$ with $v(0)=-a u(0)$, because $B$ lies strictly
between $A$ and $M$. The turn determinant at $B=\mu_{j+1}$ satisfies
\begin{equation}\label{cusppr:A-needle}
\det(v(0),w(0))=-a\det(u(0),w(0))=-a\Delta(0).
\end{equation}
It is nonzero and retains sign $-\operatorname{sgn}\Delta$ nearby.
The other needle corner is $M$, whose sign is $\tau_j$. By (i), these
signs are opposite on the no-loop side and equal on the loop side.
In case B there is $a>0$ with $u(0)=-a v(0)$. The turn determinant at
$A=\mu_{j-1}$ is therefore
\begin{equation}\label{cusppr:B-needle}
\det(h(0),u(0))=-a\det(h(0),v(0))=-a\Delta(0).
\end{equation}
It too retains sign $-\operatorname{sgn}\Delta$, while the other corner
is $M$. Applying (i) gives (ii) also in case B.

\emph{Threads and the empty condition.}
In case A, start at the newborn visit on $E_{j-1}$ and follow the
orientation. The traversal runs to $M$, along $E_j$ to $B$, and then
along $E_{j+1}$ to the other newborn visit. This is the arc in (iii).
In case B it runs from the newborn visit on $E_{j-2}$ to $A$, along
$E_{j-1}$ to $M$, and then along $E_j$ to the other visit.
Each description contains exactly the two specified polygon vertices.

No crossing other than the newborn can have both visits on this short
arc. The arc is contained in three consecutive edges. The two adjacent
pairs have no interior intersection because their incident directions
are linearly independent on a generic side. The remaining pair is the
newborn pair; its two nonparallel supporting lines meet only once, at
the newborn crossing itself. Thus any other crossing visit in the open
short arc has its partner in the complementary open arc.

If the newborn visits are cyclically adjacent in the Gauss word, at
least one of those two open arcs contains no crossing visit. A crossing
visit in the short arc would force a partner in the complementary arc,
so would make both nonempty. Hence the short arc is crossing-free.
In particular the complete edge between $c_1,c_2$ has no crossing on the
loop side. This middle edge belongs to neither member of the newborn
pair. By (i), every crossing pair involving it is unchanged, so it has
no crossing on the no-loop side either. This proves (iii) without
assuming that a general cusp is empty.

\emph{The rotation sign.}
Let $\vartheta_i(t)\in(-\pi,\pi)$ be the principal turn at vertex $i$
on a punctured side. The quotient of consecutive unit edge directions
is $\exp(\mathrm i\vartheta_i)$; multiplying around the polygon gives
$\exp(\mathrm i\sum_i\vartheta_i)=1$. Therefore
$\sum_i\vartheta_i$ is an integer multiple of $2\pi$, and
$\operatorname{rot}(P)=(2\pi)^{-1}\sum_i\vartheta_i$ is an integer.
Every principal turn is continuous away from an antiparallel pair of
incident edges. It follows that rotation is constant on each of the two
connected punctured sides.

At the centre, every turn determinant except that at $j$ is nonzero:
its three vertex indices would otherwise give a second collinear triple.
The corresponding principal turns therefore have the same limits from
the two sides. At $j$ the nonzero incident vectors are antiparallel.
The principal turn tends to $+\pi$ on the side with $T>0$ and to
$-\pi$ on the side with $T<0$. Set $s=\tau_j(P_{\rm loop})$.
Then the loop-side limit of $\vartheta_j$ is $s\pi$ and the no-loop-side
limit is $-s\pi$. Cancelling the common limits of all other turns gives
\begin{equation}\label{cusppr:rotation-limits}
2\pi\bigl(\operatorname{rot}(P_{\rm loop})-
                 \operatorname{rot}(P_{\rm no})\bigr)
=s\pi-(-s\pi)=2s\pi.
\end{equation}
Dividing by $2\pi$ proves (iv). No rotation is assigned to the singular
centre, and no differentiability or unthreadedness was used.
\end{proof}

\begin{definition}[deletion and halves]\label{def:deletion-halves}
For $n\geq4$ and $j\in\ZZ/n$, the \emph{deletion} $P\setminus j$ is the
$(n-1)$-tuple obtained by omitting $\mu_j$, with the induced cyclic order.
At a simple vertex--edge wall at $(M;a)$ with $b=a+1$, the \emph{halves} at
the centre $P^0=P(0)$ are the polygons
\[
\lambda_1=(\mu_M,\mu_{M+1},\ldots,\mu_a),\qquad
\lambda_2=(\mu_M,\mu_b,\mu_{b+1},\ldots,\mu_{M-1}),
\]
all vertices taken at $t=0$: $\lambda_1$ follows the edges $E_M,\ldots,E_{a-1}$
of $P^0$ and closes with the segment $[\mu_a,\mu_M]\subset E_a$; $\lambda_2$
opens with $[\mu_M,\mu_b]\subset E_a$, follows $E_b,\ldots,E_{M-2}$ and closes
with $E_{M-1}$.
\end{definition}

\begin{lemma}[deletions and halves are generic]\label{lem:children}
\begin{enumerate}
\item[(i)] At a simple flat wall at $j$ with $n\geq4$, $P(0)\setminus j$
is generic, and for small $t\ne0$ the deletions $P(t)\setminus j$ lie in
the chamber of $P(0)\setminus j$.
\item[(ii)] At a simple vertex--edge wall at $(M;a)$, both halves of
Definition~\ref{def:deletion-halves} are generic and have between $3$ and
$n-2$ vertices. The central parent is automatically regular.
\end{enumerate}
\status{proved (refereed: bench A, 2026-09-05T18:22:57Z; transcribed from CV lem:flatdata(ii), thm:main(H); RC flat:s3, s7:halves (C015); the two locators use different index conventions --- CV one-based tail, identical to this document; RC zero-based head with a one-step shift, its vertex $r$ being vertex $r+1$ here and its incoming edge $e_r$ the edge $E_r$ (Section~\ref{sec:conventions}) --- and the deletion $P(0)\setminus j$ and the vertex counts of the halves are stated in this document's labels (F-25-104); round~5: clause (i)'s chamber conclusion is also CV thm:main(A)(i), stated there for CV's simple flat events under its guarded genericity (Convention conv:events), whereas this document asserts it for the simple flat wall of Definition~\ref{def:walls} under the all-triple genericity of Definition~\ref{def:generic} --- the printed proof (Lemma~\ref{lem:flat-sides}(iii) with the openness of the generic locus) carries the wider scope, F-25-104 (codex B-Q1-CHILD-AI))}
\end{lemma}
\begin{proof}
Clause (i) is the deletion assertion in
Lemma~\ref{lem:flat-sides}(iii), which includes the fused-edge
genericity proof. Its chamber conclusion also follows directly from openness
of the generic locus and continuity of deletion: after shrinking, the image
of the whole parameter interval, including zero, is connected and generic,
so it lies in the connected component containing the central deletion.

For (ii), put $b=a+1$. By Lemma~\ref{lem:wall-sides} the central parent
has nonzero edges and nonzero turns. Relabel cyclically for the count only,
so $M=1$ and $a=k$. The imposed restriction
$M\notin\{a-1,a,a+1,a+2\}$ excludes $k=1,2,n-1,n$. Thus
$3\leq k\leq n-2$. The first half has $k$ vertices and the second has
$n-k+1$, which satisfies the same bounds. This count uses the stated SM1
remoteness condition; it does not add a separate regular-path hypothesis.

Each child vertex is a distinct inherited parent vertex. The only zero
parent point triple is $\{a,b,M\}$. The first half contains $a,M$ but not
$b$, and the second contains $M,b$ but not $a$, so neither contains this
triple. Every child point determinant is therefore a nonzero parent
determinant, proving (G1) for both halves.

Every inherited child edge is a parent edge. The one cut edge in either
half is a positive subsegment of $E_a$. More explicitly, for
$M=A+\lambda(B-A)$ with $0<\lambda<1$, its two possible directed vectors
are $\lambda d_a$ and $(1-\lambda)d_a$, respectively. A point in the
relative interior of either cut edge is in the relative interior of $E_a$.
If three distinct child edge interiors concurred, their parent edge images
would be three distinct parent edges with an interior concurrence: each
half has only one cut edge, so this edge map is injective. The parent edges
must be pairwise remote, because adjacent central parent edges with nonzero
turn meet only at their shared endpoint. Such a concurrence contradicts
$Z_{\rm c}=\varnothing$. Hence both halves satisfy (G2), and together with
their already proved (G1) this is exactly SM1 genericity.
\end{proof}

\subsection{Relative general position in the regular locus}

% Printed in round 1 from the Codex pack candidate_transport_v1 (C021); the statement is RC tr:factors restricted to the Delta and T classes with RC tr:bezout for the per-coordinate certificates; the coordinate-affine specialization and the irreducibility argument are printed here (F-25-183).
\begin{lemma}[the polynomial controls for SM genericity]
\label{lem:transport-polynomials}
Write the scalar vertex coordinates as $x_1,\ldots,x_{2n}$.
For each unordered triple of distinct vertex indices use one ordered
representative of
\begin{equation}\label{transport:delta}
 \Delta_{ijk}=\det(\mu_j-\mu_i,\mu_k-\mu_i).
\end{equation}
For each unordered triple of pairwise remote edges use one ordered
representative of
\begin{equation}\label{transport:T}
 T_{efg}=\det\begin{pmatrix}\ell_e\\\ell_f\\\ell_g\end{pmatrix},
 \qquad \ell_i=(-d_{i,y},d_{i,x},\det(\mu_i,d_i)),
 \quad d_i=\mu_{i+1}-\mu_i.
\end{equation}
Let $\mathcal F$ be this finite family. Every member is nonzero and
irreducible over $\mathbb R$, and distinct named members are nonassociate.
Each is affine in every scalar coordinate separately. If $F$ depends on
$x_j$, write $F=a x_j+b$; then $a$ is a nonzero polynomial in the other
coordinates. If distinct $F=a x_j+b$ and $G=c x_j+d$ both depend on $x_j$,
the polynomial $ad-bc$ is nonzero. At a specialization with $a\ne0$,
the root of $F$ in $x_j$ is simple; at a specialization with $ad-bc\ne0$,
the two polynomials have no common root.
\status{proved (refereed: bench A, 2026-09-05T15:51:04Z; transcribed from RC tr:factors restricted to the $\Delta$ and $T$ classes, with RC tr:bezout for the per-coordinate certificates (C021); the coordinate-affine specialization and the irreducibility argument are printed here (F-25-183))}
\end{lemma}
\begin{proof}
The quadratic $q(u,v)=u_xv_y-u_yv_x$ has symmetric matrix, in the order
$(u_x,u_y,v_x,v_y)$,
\begin{equation}\label{transport:qmatrix}
 \frac12\begin{pmatrix}0&0&0&1\\0&0&-1&0\\0&-1&0&0\\1&0&0&0\end{pmatrix}.
\end{equation}
Its rank is four. A nontrivial factorization of a homogeneous quadratic is
a product of two homogeneous linear forms. If their coefficient vectors are
$a,b$, the symmetric matrix of their product is
$(ab^{\mathsf T}+ba^{\mathsf T})/2$, of rank at most two. Thus $q$ is
irreducible. The vector pairs defining $\Delta$ and the auxiliary remote
direction determinant $H_{ef}=\det(d_e,d_f)$ are independent linear coordinates:
complete either surjective linear map to an invertible coordinate change.
The polynomial is then $q$ with extra unused variables. A factorization cannot
depend on an unused variable because polynomial degrees in that variable add.
This proves irreducibility and nonzeroness of $\Delta$ and $H$.

For $T$, the six endpoint occurrences are independent. Use the invertible
tail/direction coordinates $(\mu_e,d_e),(\mu_f,d_f),(\mu_g,d_g)$.
Writing the cross product of the first two line rows as $(A,B,C)$ gives
$C=H_{ef}$. Expansion in the third row yields
\begin{equation}\label{transport:T-expand}
 T=d_{g,y}(\mu_{g,x}C-A)+d_{g,x}(B-\mu_{g,y}C).
\end{equation}
These two coefficients have no common nonconstant divisor. Indeed any
irreducible common divisor is independent of $\mu_{g,x}$ because the second
coefficient is, and independent of $\mu_{g,y}$ because the first is.
Comparing coefficients would make it divide $A,B,C$. But $C=H_{ef}$ is
irreducible and does not divide both $A,B$: take $d_e=d_f=(1,0)$ and tails on
different horizontal lines. Then $C=0$ while $A\ne0$. Therefore the two
coefficients in \eqref{transport:T-expand} are coprime. In a factorization
of a polynomial homogeneous of degree one in $(d_{g,x},d_{g,y})$, one factor
is independent of both variables and divides both coefficients. It must be
a unit. This proves irreducibility of $T$; the same expansion proves it is
nonzero.

Each $\Delta$ depends on every vertex in its three-element support. Each
$T$ depends on every endpoint in its six-element support: with one endpoint
fixed, the line through it can be varied by moving its other endpoint;
choose the other two lines to intersect away from the fixed endpoint to
make the determinant vary. Thus different supports cannot give associates,
and the three- and six-vertex families cannot give associates. For a fixed
six-vertex support, the three pairwise remote cycle edges form a perfect
matching. A proper subset of a cycle induces disjoint paths, each of which
has at most one perfect matching (remove its first edge recursively).
The only possible two matchings occur when $n=6$ and the support is the
whole cycle; they are the two alternating matchings. At the specialization
\begin{equation}\label{transport:T-n6-witness}
 (\mu_1,\ldots,\mu_6)=((1,0),(2,0),(0,1),(0,2),(1,1),(2,2))
\end{equation}
the tail-indexed values are $T_{135}=0$ and $T_{246}=2$. Hence they too are
nonassociate. One representative per unordered point triple already removes
the sign associates among the $\Delta$'s.

Every scalar endpoint coordinate occurs affinely in its edge line row, and
the three rows of $T$ have disjoint endpoint supports. Hence $T$ is affine
in each scalar coordinate. The same assertion for $\Delta$ follows by
expanding its area determinant. Genuine dependence therefore has a nonzero
slope $a$. Irreducibility of $a x_j+b$ implies that $a,b$ are coprime in
the polynomial ring of the other coordinates: a nonconstant common divisor
would factor $F$. Likewise $c,d$ are coprime. If $ad-bc=0$, coprimality
of $a,b$ in this unique factorization domain implies $a\mid c$; put $c=au$.
The identity then gives $d=bu$. Coprimality of $c,d$ forces $u$ to be a unit,
so $G=uF$, contradicting nonassociation. Thus $ad-bc$ is a nonzero polynomial.
After specialization, a nonzero slope makes a linear root simple, and
$cF-aG=cb-ad$ excludes a common root when $ad-bc\ne0$.
\end{proof}

\begin{theorem}[relative general position and the wall dictionary]
\label{thm:relgp}
Let $\gamma:[0,1]\to\mathcal R_n$ be continuous with generic endpoints,
and let $\delta>0$. There is a piecewise-affine path in $\mathcal R_n$ with
the same labelled endpoints, uniformly within $\delta$ of $\gamma$, which
is collision-free and whose nongeneric parameters are finitely many isolated
simple wall germs, each of exactly one of the types (F), (V) bigon, (V)
sliding, (T), (E), (C). No cusp occurs. Every occurrence label is retained.
\status{proved (refereed: bench A, 2026-09-05T15:51:04Z and 2026-09-06T01:50:42Z; transcribed from RC tr:gp, tr:five-events (C021))}
\end{theorem}
\begin{proof}
\emph{Open neighbourhoods and endpoint collars.}
The regular locus is open: nonzero edge lengths and avoidance of antiparallel
direction pairs are open conditions. The generic locus is open as well.
To see this directly under SM (G1),(G2), finitely many nonzero point determinants
retain their signs near a generic tuple. Every central remote pair is either
compact-disjoint or an interior transverse pair (Lemma~\ref{lem:g1}(v));
disjointness persists, and a transverse interior crossing persists with
continuous parameters, by Lemma~\ref{lem:wall-segment-stability}. Distinct crossing points have positive pairwise separation, since
a coincidence would place at least three edge interiors at one point. Thus
no triple interior concurrence appears in a sufficiently small neighbourhood.
This proves openness even when a concurrence polynomial vanishes at an
inactive configuration, such as three line extensions concurrent outside the
segments. No all-polynomials-nonzero endpoint assumption is needed.

Cover the compact image of $\gamma$ by axis-parallel open cubes with closures
in $\mathcal R_n$ and diameters less than $\delta$. The endpoint cubes may
be chosen inside the generic locus. A sufficiently fine subdivision
$0=t_0<\cdots<t_N=1$ has each image
$\gamma([t_{i-1},t_i])$ inside one such cube $C_i$, with the first and last
cubes generic. Consecutive cubes overlap in a nonempty open set because they
both contain $\gamma(t_i)$. Refine if necessary so $N\geq3$.

Choose independent internal waypoints $z_i\in C_i\cap C_{i+1}$, while
$z_0=\gamma(0)$ and $z_N=\gamma(1)$ remain fixed. Connect $z_{i-1}$ to $z_i$
by changing their $2n$ scalar coordinates one at a time in a fixed order.
Every intermediate hybrid belongs to $C_i$, since an axis-parallel cube is
closed under coordinate hybridization. Parameterize these legs over
$[t_{i-1},t_i]$. The old and new points at each time belong to the same
cube, proving the uniform $\delta$ bound. All legs are regular. The first
and last coordinate polylines are entirely generic; they are endpoint collars.
Only legs joining two independently variable internal waypoints need further
conditions.

\emph{One joint algebraic choice for the central legs.}
The product of the open waypoint sets is an open subset of a finite-dimensional
real coordinate space. For each central leg moving scalar coordinate $x_j$,
impose the following nonvanishing requirements:
\begin{enumerate}
\item Every $F\in\mathcal F$ is nonzero at both endpoint hybrids.
\item For every $x_j$-dependent $F=a x_j+b$, the fixed-coordinate slope $a$
is nonzero; for every distinct such pair $F,G$, its fixed-coordinate
linear resultant $ad-bc$ is nonzero.
\item The scalar step is nonzero, and both coordinate differences of every
pair of vertex occurrences are nonzero at the endpoint hybrids.
\end{enumerate}
Every requirement is the nonvanishing of a nonzero polynomial in the joint
waypoint variables. For the first two assertions, a central hybrid selects
each coordinate from one of two independent internal waypoints; its projection
contains an open coordinate box. A nonzero polynomial in the hybrid coordinates,
or in the fixed coordinates of the leg, therefore cannot pull back to the zero
polynomial. The nonzero slopes and resultants are supplied by
Lemma~\ref{lem:transport-polynomials}. Coordinate-difference polynomials use
distinct occurrence variables; the scalar step uses two independent values
of the moving coordinate, so these polynomials also are nonzero.

There are finitely many conditions. Their product is a nonzero polynomial,
which cannot vanish on a nonempty real open set. For completeness, this last
fact follows by induction on the number of variables: a univariate nonzero
polynomial has finitely many roots; on a product of open intervals, fixing all
but one variable would force every coefficient polynomial to vanish, and the
induction hypothesis then forces the original polynomial to be zero.
Thus one joint waypoint choice satisfies all requirements. Fixed original
endpoints are excluded from this choice, which is why collars were separated.

\emph{Finite simple roots and collision avoidance.}
On a central leg a factor independent of $x_j$ stays at its nonzero endpoint
value. A dependent factor is a genuinely linear polynomial in the moving
coordinate with nonzero slope, hence has at most one root on the leg.
Different factors have no common root by their nonzero linear resultant.
No root is at a leg endpoint, and each root changes the sign of exactly one
named factor because the scalar step is nonzero. Consequently all central
roots are finite, isolated, smooth hypersurface crossings. Tangencies and
intersections of distinct control hypersurfaces have been avoided by the
explicit slope and resultant conditions, not by an unproved generic-position
assertion.

On a leg moving the $x$-coordinate of one vertex, its $y$-difference from
every other vertex is fixed and nonzero. It cannot collide with any of them.
All other vertex pairs are fixed and distinct. The same argument exchanges
$x,y$ on a $y$-coordinate leg. The collars are generic and hence have distinct
vertices. Thus the whole resulting path is collision-free, including its
fixed endpoints. Every physical edge occurrence stays nonzero because all
legs lie in their regular cubes. A chosen labelled edge can therefore be
followed along the path without relabelling or deletion; no root-independence
statement is used.

If all members of $\mathcal F$ are nonzero, (G1) holds. Any triple interior
concurrence would, under (G1), use pairwise remote edges and make one $T$
zero, so (G2) holds too. Hence every nongeneric parameter is among the finite
central roots. Some roots of an inactive $T$ may still be generic; they are
not declared wall events. Collars contain no nongeneric parameters even when
an inactive polynomial there has nonisolated zeros.

\emph{A point-determinant root.}
At a root of $\Delta_K$, exactly the point triple $K$ is collinear, all
vertices remain distinct, and all $T$'s are nonzero. Thus $Z_{\rm pt}=\{K\}$
and $Z_{\rm c}=\varnothing$. The determinant changes sign simply on the leg.
Count the cyclically adjacent pairs in $K$. If there are at least two and
$n\geq4$, $K=\{j-1,j,j+1\}$ is a consecutive triple. Its middle vertex
must lie strictly between its neighbours, because otherwise the two consecutive
directions there are antiparallel, contrary to regularity. This is exactly (F).
When $n=3$, a distinct collinear triple forms a closed collinear triangle and
necessarily has an antiparallel consecutive pair; such a root cannot lie in
the regular cube.

If there is exactly one adjacent pair, write it as $a,b=a+1$ and call the
third index $M$. The lack of a second adjacent pair is precisely
$M\notin\{a-1,a,a+1,a+2\}$. The contact point is not either endpoint, by
collision freedom. If $M$ is strictly inside $[\mu_a,\mu_b]$, this is (V).
Both neighbours of $M$ have nonzero heights over that line, since another
zero height would be another point triple. Equal height signs give the bigon
branch and opposite signs the sliding branch. If $M$ lies outside the closed
segment, this is (E). Finally, no adjacent pair gives exactly the pairwise
nonadjacent point triple of (C). These alternatives are exhaustive and disjoint.

\emph{An active concurrence root.}
At a root of a single $T_{efg}$ all point determinants are nonzero. If it is
nongeneric, some three edge interiors concur; they must be pairwise remote,
and their concurrence polynomial must be this unique zero member. Thus
$Z_{\rm pt}=\varnothing$ and $Z_{\rm c}=\{\{e,f,g\}\}$. At their common
interior point every pair is transverse, since coincident remote supporting
lines would force a point-determinant zero. All three crossing pairs persist
on nearby punctured sides by their continuous interior intersection parameters.

To verify the sign-change part of (T), put $H_{ef}=\det(d_e,d_f)$ and write
$q_{ef}=\mu_e+t_{ef}d_e$. The cross product of the first two line rows in
\eqref{transport:T} equals $H_{ef}(q_{ef,x},q_{ef,y},1)$, as follows from
its third coordinate $H_{ef}$ and the two line equations. Therefore
\begin{equation}\label{transport:T-order}
 T_{efg}=H_{ef}\det(d_g,q_{ef}-\mu_g)
        =-(t_{ef}-t_{eg})H_{ef}H_{eg}.
\end{equation}
The second equality subtracts $q_{eg}=\mu_e+t_{eg}d_e$, which lies on
the $g$ line, and uses $\det(d_g,d_e)=-H_{eg}$. The two $H$ factors are
nonzero with constant sign near the root. The simple sign change of $T$
therefore changes the sign of $t_{ef}-t_{eg}$. Apply the same calculation
after cyclic permutation of $(e,f,g)$ to obtain the other two required
parameter-difference sign changes. This is exactly (T).

Every nongeneric parameter has now been classified as one named simple wall.
There are no cusp walls, because their central antiparallel pair is excluded
throughout the regular cubes. Both the endpoints and every occurrence label
have been kept exactly as required.
\end{proof}

